\documentclass[11pt]{amsart}

\usepackage[utf8]{inputenc}
\usepackage{amsmath,amssymb,amsthm}
\usepackage[hidelinks]{hyperref}

\theoremstyle{plain}
\newtheorem{theorem}{Theorem}
\newtheorem{lemma}[theorem]{Lemma}
\newtheorem{proposition}[theorem]{Proposition}
\newtheorem{conjecture}[theorem]{Conjecture}
\newtheorem{claim}[theorem]{Claim}
\newtheorem{question}[theorem]{Question}

\theoremstyle{definition}
\newtheorem{remark}[theorem]{Remark}

\def\bE{\mathbb{E}}
\def\bR{\mathbb{R}}
\def\Pb{\mathbb{P}}

\begin{document}

\title{On the Probability a Weighted Bernoulli Sum Exceeds Its Mean}
\author{}
\date{}
\maketitle

\section{Statement}

Let $w\in\bR_{>0}^m$ be a vector of positive weights with $\sum_{i=1}^m w_i=1$, and
let $v\in\{0,1\}^m$ be a random binary vector whose coordinates are i.i.d.\
$\mathrm{Bernoulli}(p)$. Set $X=\langle w,v\rangle$. Then $X$ takes values in $[0,1]$
and $\bE[X]=p$. We study the following question.

\begin{question}\label{q:prob}
For which $p\in(0,1)$ does the following hold? For every $w\in\bR_{>0}^m$ with
$\sum_{i=1}^m w_i=1$, if $v\in\{0,1\}^m$ has i.i.d.\ $\mathrm{Bernoulli}(p)$
coordinates, then
\[
\Pb[\langle v,w\rangle\ge p]\ge p .
\]
\end{question}

We relate Question~\ref{q:prob} to the Manickam--Mikl\'os--Singhi (MMS) conjecture.

\begin{conjecture}[MMS]\label{conj:MMS}
Let $n,k$ be positive integers with $n\ge 4k$. For any real numbers
$x_1,\dots,x_n$ with $\sum_{i=1}^n x_i=0$, at least $\binom{n-1}{k-1}$ of the
$k$-element subsets $S\subseteq[n]$ satisfy $\sum_{i\in S}x_i\ge 0$.
\end{conjecture}

A linear bound (the conjecture holds for all $n\ge ck$ with an absolute constant
$c$, e.g.\ $c=10^{46}$) is known by Pokrovskiy. The aim of this note is the
following conditional theorem.

\begin{theorem}\label{thm:main}
Suppose there is a constant $c>0$ such that Conjecture~\ref{conj:MMS} holds for
all positive integers $n,k$ with $n\ge ck$. Then Question~\ref{q:prob} holds for
all $p\in[0,1/c)$.
\end{theorem}

In what follows we only ever apply Conjecture~\ref{conj:MMS} through its hypotheses
$\sum_i x_i=0$ and $n\ge ck$; no structural restriction on the values $x_i$ is
used. (In Remark~\ref{rem:weakMMS} we observe, separately and without affecting the
main argument, that a structurally weaker form of the conjecture would already
suffice.)

\section{Proof of Theorem~\ref{thm:main}}

The case $p=0$ is trivial (the inequality $\Pb[X\ge 0]\ge 0$ always holds), so
assume $p>0$. We first prove the statement for rational $p\in(0,1/c)$, and then
extend it to all real $p\in(0,1/c)$ by a limiting argument.

\subsection{Two elementary facts}

We isolate two standard facts used below.

\begin{lemma}[Left-continuity of the upper tail]\label{lem:leftcont}
Let $Z$ be a real-valued random variable and define
$g(s)=\Pb[Z\ge s]$ for $s\in\bR$. Then $g$ is left-continuous: if
$s_1\le s_2\le\cdots$ and $s_i\to s$, then $g(s)=\lim_{i\to\infty}g(s_i)$.
\end{lemma}

\begin{proof}
Since $s_i\le s$, the event $\{Z\ge s\}$ is contained in $\{Z\ge s_i\}$ for every
$i$, and the events $\{Z\ge s_i\}$ are nested decreasing in $i$ (as $s_i$ is
nondecreasing). We claim
\[
\{Z\ge s\}=\bigcap_{i\ge1}\{Z\ge s_i\}.
\]
The inclusion ``$\subseteq$'' is immediate. For ``$\supseteq$'', if $Z\ge s_i$ for
all $i$, then $Z\ge\sup_i s_i=\lim_i s_i=s$. By continuity of probability from
above for a decreasing sequence of events,
$g(s)=\Pb\big[\bigcap_i\{Z\ge s_i\}\big]=\lim_i\Pb[Z\ge s_i]=\lim_i g(s_i)$.
\end{proof}

\begin{lemma}[Monotone coupling]\label{lem:coupling}
Fix $w\in\bR_{>0}^m$ and $0<q\le p<1$. Let $v$ have i.i.d.\
$\mathrm{Bernoulli}(p)$ coordinates and let $v'$ have i.i.d.\
$\mathrm{Bernoulli}(q)$ coordinates. Then for every threshold $s\in\bR$,
\[
\Pb[\langle w,v\rangle\ge s]\ \ge\ \Pb[\langle w,v'\rangle\ge s].
\]
\end{lemma}

\begin{proof}
Let $U_1,\dots,U_m$ be i.i.d.\ uniform on $[0,1]$ and set
$v_j=\mathbf 1[U_j\le p]$ and $v'_j=\mathbf 1[U_j\le q]$. Then $v$ has i.i.d.\
$\mathrm{Bernoulli}(p)$ coordinates, $v'$ has i.i.d.\
$\mathrm{Bernoulli}(q)$ coordinates, and since $q\le p$ we have $v'_j\le v_j$ for
every $j$ pointwise. As all $w_j>0$,
$\langle w,v\rangle\ge\langle w,v'\rangle$ pointwise on this probability space.
Hence $\{\langle w,v'\rangle\ge s\}\subseteq\{\langle w,v\rangle\ge s\}$, which
gives the claimed inequality.
\end{proof}

\subsection{The rational case}

Fix a rational $p=k_0/n_0\in(0,1/c)$ with positive integers $k_0,n_0$, and fix a
weight vector $w\in\bR_{>0}^m$ with $\sum_{j=1}^m w_j=1$. Let $v$ have i.i.d.\
$\mathrm{Bernoulli}(p)$ coordinates and $X=\langle w,v\rangle$. We must show
$\Pb[X\ge p]\ge p$.

\paragraph{A combinatorial model.}
Let $t\ge1$ be an integer parameter and set $k=k_0t$ and $n=n_0t$, so that
$k/n=k_0/n_0=p$ for every $t$. (For the main argument any single choice of $t$,
e.g.\ $t=1$, would do; we keep $t$ free only so that the same model can be reused
in Remark~\ref{rem:weakMMS}.) Construct a random $0/1$ matrix $M$ with
$n$ rows and $m$ columns as follows. Let $C_1,\dots,C_m$ be i.i.d.\ random
variables, each uniform on $[n]=\{1,\dots,n\}$; here $C_j$ records the row in which
column $j$ places its unique $1$. Formally, the $(i,j)$ entry of $M$ is
$M_{ij}=\mathbf 1[C_j=i]$, so column $j$ has a single $1$ (in row $C_j$) and $0$'s
elsewhere. Denote the rows of $M$ by $r_1,\dots,r_n\in\{0,1\}^m$, where the
$j$-th coordinate of $r_i$ is $(r_i)_j=M_{ij}=\mathbf 1[C_j=i]$, and put
$y_i=\langle w,r_i\rangle$ for $i\in[n]$. Independently of $M$ (equivalently, of
$C_1,\dots,C_m$), choose a uniformly random $k$-element subset $S\subseteq[n]$, and
set $Y=\sum_{i\in S}y_i$.

\begin{claim}\label{cl:dist}
$Y$ has the same distribution as $X=\langle w,v\rangle$, for every choice of the
parameter $t\ge1$.
\end{claim}

\begin{proof}
We have $Y=\big\langle w,\ u\big\rangle$ with $u:=\sum_{i\in S}r_i\in\{0,1\}^m$, so
it suffices to show that $u$ has the same distribution as $v$, i.e.\ that $u$ has
i.i.d.\ $\mathrm{Bernoulli}(p)$ coordinates.

\medskip
\emph{Step 1: the conditional law of $u$ given $S$.}
Fix an arbitrary $k$-element set $S\subseteq[n]$. Conditionally on $S$, the column
positions $C_1,\dots,C_m$ retain their joint law (they are independent of the
choice of $S$), namely they are i.i.d.\ uniform on $[n]$. For each $j\in[m]$, the
$j$-th coordinate of $u$ is
\[
u_j=\sum_{i\in S}(r_i)_j=\sum_{i\in S}\mathbf 1[C_j=i]=\mathbf 1[C_j\in S],
\]
where the last equality holds because $C_j$ takes exactly one value, so at most one
indicator $\mathbf 1[C_j=i]$ (with $i\in S$) is nonzero. Thus
$u_j=\mathbf 1[C_j\in S]$ is a measurable function of $C_j$ alone.

Since $C_1,\dots,C_m$ are independent, the random variables
$u_1=f(C_1),\dots,u_m=f(C_m)$ with $f(c):=\mathbf 1[c\in S]$ are independent as
well (functions of independent random variables are independent). Moreover, for
each $j$,
\[
\Pb[u_j=1\mid S]=\Pb[C_j\in S\mid S]=\frac{|S|}{n}=\frac{k}{n}=p,
\]
because $C_j$ is uniform on the $n$ rows and exactly $|S|=k$ of them lie in $S$.
Hence, conditionally on $S$, the coordinates $u_1,\dots,u_m$ are
\emph{independent} and each is $\mathrm{Bernoulli}(p)$; equivalently, for every
subset $T\subseteq[m]$,
\[
\Pb\big[u_j=1\ \text{for all } j\in T\ \big|\ S\big]
=\prod_{j\in T}\Pb[C_j\in S\mid S]=\Big(\frac{k}{n}\Big)^{|T|}=p^{|T|}.
\]
Consequently, for every fixed $b=(b_1,\dots,b_m)\in\{0,1\}^m$,
\begin{equation}\label{eq:condlaw}
\Pb[u=b\mid S]=\prod_{j=1}^m p^{\,b_j}(1-p)^{\,1-b_j},
\end{equation}
which is precisely the probability that $v=b$. In particular, this conditional law
does \emph{not} depend on the choice of $S$ (it only used $|S|=k$ and the
independence/uniformity of $C_1,\dots,C_m$). This is the point the marginal
computation alone misses: equal one-coordinate marginals would not suffice, but
\eqref{eq:condlaw} pins down the full joint law.

\medskip
\emph{Step 2: averaging over $S$.}
By the law of total probability, since $S$ is independent of
$(C_1,\dots,C_m)$, for every $b\in\{0,1\}^m$,
\[
\Pb[u=b]=\sum_{|S|=k}\Pb[S\ \text{chosen}]\;\Pb[u=b\mid S]
=\Big(\sum_{|S|=k}\Pb[S\ \text{chosen}]\Big)\prod_{j=1}^m p^{\,b_j}(1-p)^{\,1-b_j},
\]
where we used that the conditional law \eqref{eq:condlaw} is the same constant
(in $S$) for every $k$-set $S$. Since $\sum_{|S|=k}\Pb[S\ \text{chosen}]=1$, we
obtain
\[
\Pb[u=b]=\prod_{j=1}^m p^{\,b_j}(1-p)^{\,1-b_j}=\Pb[v=b]
\qquad\text{for all } b\in\{0,1\}^m.
\]
Therefore $u$ and $v$ have the same distribution on $\{0,1\}^m$, and hence
$Y=\langle w,u\rangle$ has the same distribution as $X=\langle w,v\rangle$. Note
that this conclusion is independent of $t$: only the identity $k/n=p$ was used.
\end{proof}

\begin{claim}\label{cl:bound}
$\Pb[Y\ge p]\ge p$.
\end{claim}

\begin{proof}
Define $z_i=y_i-\frac1n$ for $i\in[n]$. Since each column of $M$ contains exactly
one $1$, we have $\sum_{i=1}^n r_i=\mathbf 1$ (the all-ones vector in
$\bR^m$), hence
\[
\sum_{i=1}^n y_i=\Big\langle w,\sum_{i=1}^n r_i\Big\rangle
=\langle w,\mathbf 1\rangle=\sum_{j=1}^m w_j=1,
\]
and therefore
\begin{equation}\label{eq:sumzero}
\sum_{i=1}^n z_i=\sum_{i=1}^n y_i-n\cdot\frac1n=0 .
\end{equation}
Note that \eqref{eq:sumzero} holds for \emph{every} realization of $M$.

Since $p=k/n<1/c$, we have $n>ck$, hence $n\ge ck$. Condition on any realization
of $M$; this fixes the real numbers $z_1,\dots,z_n$, which satisfy
\eqref{eq:sumzero}. Applying the hypothesis (Conjecture~\ref{conj:MMS} for
$n\ge ck$) to $z_1,\dots,z_n$ --- with \emph{no} restriction on the values
$z_i$, which may be arbitrary reals --- at least $\binom{n-1}{k-1}$ of the
$k$-element subsets $S\subseteq[n]$ satisfy $\sum_{i\in S}z_i\ge0$. Because $S$ is
chosen uniformly among all $\binom{n}{k}$ such subsets, independently of $M$,
\[
\Pb\Big[\textstyle\sum_{i\in S}z_i\ge0 \ \Big|\ M\Big]
\ \ge\ \frac{\binom{n-1}{k-1}}{\binom{n}{k}}
=\frac{k}{n}=p ,
\]
and this bound holds for every realization of $M$. Taking expectation over $M$,
\begin{equation}\label{eq:avg}
\Pb\Big[\textstyle\sum_{i\in S}z_i\ge0\Big]\ \ge\ p .
\end{equation}
Finally, $\sum_{i\in S}z_i=\sum_{i\in S}y_i-\frac{|S|}{n}=Y-\frac{k}{n}=Y-p$, so
the event $\{\sum_{i\in S}z_i\ge0\}$ coincides with $\{Y\ge p\}$. Substituting
into \eqref{eq:avg} gives $\Pb[Y\ge p]\ge p$.
\end{proof}

Combining Claims~\ref{cl:dist} and~\ref{cl:bound},
\[
\Pb[X\ge p]=\Pb[Y\ge p]\ge p ,
\]
which proves Question~\ref{q:prob} for every rational $p\in(0,1/c)$.

\subsection{A structural remark: a weaker MMS already suffices}

We now record an observation about the values $z_i$ that appear in the model. It is
not used anywhere in the proof of Theorem~\ref{thm:main} above --- that proof
applies the \emph{full} Conjecture~\ref{conj:MMS} to the (arbitrary real) numbers
$z_1,\dots,z_n$. The point of this subsection is only to show that the conclusion of
the theorem would already follow from a structurally weaker hypothesis, in which
all negative entries of the sequence are required to be equal.

We first determine the values taken by $y_i$.

\begin{lemma}[Sub-lemma 2a]\label{lem:values}
In the model above, for each $i\in[n]$,
\[
y_i=\sum_{j:\,M_{ij}=1} w_j=\sum_{j:\,C_j=i} w_j .
\]
Consequently either $y_i=0$ (when no column places its $1$ in row $i$, i.e.\
$C_j\ne i$ for all $j$), or $y_i$ is a sum of one or more of the weights $w_j$ and
hence
\[
y_i\ \ge\ \min_{j\in[m]} w_j\ =:\ \mu\ >\ 0 .
\]
In all cases $y_i\in\{0\}\cup[\mu,\infty)$, and this holds for every realization of
$M$.
\end{lemma}

\begin{proof}
By definition $y_i=\langle w,r_i\rangle=\sum_{j=1}^m w_j (r_i)_j$ and
$(r_i)_j=M_{ij}=\mathbf 1[C_j=i]$, so only the columns $j$ with $C_j=i$ contribute,
giving $y_i=\sum_{j:\,C_j=i} w_j$. If the index set $\{j:C_j=i\}$ is empty the sum
is $0$. Otherwise it contains at least one index $j_0$, and since all weights are
positive,
\[
y_i=\sum_{j:\,C_j=i} w_j\ \ge\ w_{j_0}\ \ge\ \min_{j\in[m]} w_j=\mu>0 .
\]
The positivity $\mu>0$ holds because $w\in\bR_{>0}^m$, and $\mu$ depends only on
$w$ (equivalently on $m$ and the fixed weights), not on $M$, $t$, or $S$.
\end{proof}

\begin{remark}[A weakened MMS would suffice]\label{rem:weakMMS}
Translating Lemma~\ref{lem:values} to $z_i=y_i-\tfrac1n$, every $z_i$ lies in
\[
\{-\tfrac1n\}\ \cup\ [\mu-\tfrac1n,\ \infty),
\]
the value $-\tfrac1n$ occurring exactly when $y_i=0$. Here $\mu=\min_j w_j>0$ is a
\emph{fixed} constant depending only on $w$, while $\tfrac1n=\tfrac1{n_0 t}\to0$ as
$t\to\infty$. Hence there is a threshold $t_0=t_0(w,n_0)$ such that for all
$t\ge t_0$ we have $\tfrac1n=\tfrac1{n_0t}<\mu$, and then every nonzero $y_i$ gives
$z_i=y_i-\tfrac1n\ge\mu-\tfrac1n>0$. For such $t$ the \emph{only} negative value
taken by any $z_i$ is $-\tfrac1n$ (attained precisely on the rows $i$ with
$y_i=0$); all other $z_i$ are strictly positive.

We may indeed take $t$ this large without changing anything: by
Claim~\ref{cl:dist} the law of $Y$ equals that of $X$ for \emph{every} $t\ge1$
(only $k/n=p$ was used), so replacing $t$ by some $t\ge t_0$ leaves the identity
$\Pb[X\ge p]=\Pb[Y\ge p]$ intact, while \eqref{eq:sumzero} still gives
$\sum_{i=1}^n z_i=0$ and $n\ge ck$ still holds (both are unaffected by $t$).

Consequently, the deduction in Claim~\ref{cl:bound} would go through verbatim if,
instead of the full Conjecture~\ref{conj:MMS}, one only assumed the following
weaker statement: \emph{for every $n\ge ck$ and all reals $x_1,\dots,x_n$ with
$\sum_i x_i=0$ in which all negative entries are equal, at least
$\binom{n-1}{k-1}$ of the $k$-subsets have nonnegative sum.} Indeed, for $t\ge t_0$
the sequence $z_1,\dots,z_n$ has exactly this structure (its only negative value is
$-\tfrac1n$). Thus Theorem~\ref{thm:main} remains valid under this weakened
hypothesis. We stress that this is an optional strengthening of the input
assumption: the proof of Theorem~\ref{thm:main} given above does not rely on it and
uses the full Conjecture~\ref{conj:MMS}.
\end{remark}

\subsection{Extension to irrational $p$}

Let $p\in(0,1/c)$ be arbitrary (in particular possibly irrational), fix
$w\in\bR_{>0}^m$ with $\sum_j w_j=1$, let $v$ have i.i.d.\
$\mathrm{Bernoulli}(p)$ coordinates, and set $X=\langle w,v\rangle$. Choose a
sequence of rationals $p_1\le p_2\le\cdots$ with $0<p_i<p$ and $p_i\to p$; this is
possible since $p>0$.

For each $i$, the rational case applied with parameter $p_i$ (and the same $w$)
gives: if $v^{(i)}$ has i.i.d.\ $\mathrm{Bernoulli}(p_i)$ coordinates, then
\begin{equation}\label{eq:rat}
\Pb[\langle w,v^{(i)}\rangle\ge p_i]\ \ge\ p_i .
\end{equation}
Since $p_i\le p$, Lemma~\ref{lem:coupling} (with $q=p_i$ and threshold $s=p_i$)
yields
\begin{equation}\label{eq:dom}
\Pb[\langle w,v\rangle\ge p_i]\ \ge\ \Pb[\langle w,v^{(i)}\rangle\ge p_i]\ \ge\ p_i ,
\end{equation}
where the last inequality is \eqref{eq:rat}. Thus $\Pb[X\ge p_i]\ge p_i$ for all
$i$.

Now apply Lemma~\ref{lem:leftcont} to the random variable $Z=X$ and the increasing
sequence $s_i=p_i\to p$:
\[
\Pb[X\ge p]=\lim_{i\to\infty}\Pb[X\ge p_i]\ \ge\ \lim_{i\to\infty}p_i=p ,
\]
where the inequality uses \eqref{eq:dom}. This proves Question~\ref{q:prob} for
$p$.

Since $p\in(0,1/c)$ was arbitrary and the case $p=0$ is trivial, the theorem
follows. \qed

\section{Optimality outside the threshold}

For completeness we record that the answer to Question~\ref{q:prob} is negative
for $p\in(1/3,1)\setminus\{1/2\}$, so a positive answer can hold only on an initial
segment together with the isolated point $p=1/2$.

\begin{proposition}\label{prop:counterexample}
The answer to Question~\ref{q:prob} is negative for every $p\in(1/3,1)$ with
$p\ne1/2$.
\end{proposition}

\begin{proof}
Suppose first $1/3<p<1/2$ and take $w=(1/3,1/3,1/3)$. Then
$\langle w,v\rangle\ge p$ iff at least two coordinates of $v$ equal $1$ (since the
attainable values of $\langle w,v\rangle$ are $0,1/3,2/3,1$ and $1/3<p<2/3$).
Hence
\[
\Pb[\langle w,v\rangle\ge p]=\binom{3}{2}p^2(1-p)+p^3=3p^2-2p^3=p(3p-2p^2).
\]
Then
\[
p-\Pb[\langle w,v\rangle\ge p]=p\big(1-3p+2p^2\big)=p(1-2p)(1-p)>0
\]
for $p\in(1/3,1/2)$, so $\Pb[\langle w,v\rangle\ge p]<p$.

If $1/2<p<1$, take $w=(1/2,1/2)$. Then $\langle w,v\rangle\ge p$ iff both
coordinates of $v$ equal $1$ (the attainable values are $0,1/2,1$ and $1/2<p<1$),
so $\Pb[\langle w,v\rangle\ge p]=p^2<p$.
\end{proof}

\end{document}
